%======================================================================
\chapter{Conclusion and Future Work}
\label{ch:conclusion}
%======================================================================

%----------------------------------------------------------------------
\section{Summary}

This project developed a modular, 2-D PDE-constrained optimal control
simulation for MRI-guided microwave thermal ablation.
The Pennes bioheat equation and Arrhenius damage ODE are discretized
by finite differences and integrated with a shared forward simulator
used by all four solver backends (open-loop, indirect PMP, direct NLP,
and receding-horizon MPC).

The indirect solver derives costate equations analytically from the
Pontryagin Maximum Principle and updates the control schedule via
forward-adjoint sweep iteration.
The direct solver delegates optimization to \texttt{scipy.optimize}
and trades speed for implementation simplicity.
The MPC solver demonstrates closed-loop robustness to simulated MRI
thermometry noise.

%----------------------------------------------------------------------
\section{Remaining Challenges}

\begin{enumerate}
    \item \textbf{MPC inner optimizer:}
        The current \texttt{\_mpc\_horizon\_optimize} uses a bang-off-bang
        heuristic rather than a true OCP solver.
        Replacing it with \texttt{solve\_direct()} or \texttt{solve\_indirect()}
        on the prediction window would yield a proper MPC.

    \item \textbf{Indirect convergence:}
        The gradient projection algorithm is not globally convergent;
        the relaxation parameter $\beta$ requires manual tuning.
        Line-search or conjugate-gradient extensions would improve robustness.

    \item \textbf{Free final time:}
        $t_f$ is currently fixed at \SI{300}{\second}.
        Extending to a true free-final-time formulation requires an
        additional transversality condition and a variable horizon.

    \item \textbf{3-D domain:}
        Clinical ablation is three-dimensional.
        The FD grid scales as $O(N^{3/2})$ in 3-D; a multigrid or
        spectral preconditioner would be needed.
\end{enumerate}

%----------------------------------------------------------------------
\section{Future Work}

\begin{itemize}
    \item Replace single-shooting direct method with direct collocation
          (e.g., CasADi/IPOPT) for better scalability.
    \item Incorporate real MRI thermometry data (proton resonance
          frequency shift maps) as the feedback signal.
    \item Extend to heterogeneous tissue properties and patient-specific
          geometry from MRI segmentation.
    \item Implement safety-critical control (control barrier functions)
          to enforce the healthy-tissue temperature constraint as a
          hard constraint rather than a soft penalty.
\end{itemize}

%----------------------------------------------------------------------
\section{Expected Contributions}

\begin{enumerate}
    \item Validated 2-D finite-difference bioheat + damage simulator
          with four interchangeable solver backends.
    \item Adjoint solver verified against finite-difference gradient checks.
    \item Quantitative comparison of open-loop, indirect, direct, and
          MPC strategies under matched conditions.
    \item Configurable boundary conditions (Dirichlet / Neumann / Robin)
          exposed as a command-line argument.
\end{enumerate}
